\documentclass[a4paper,12pt]{article}

\usepackage[T2A]{fontenc}
\usepackage[utf8]{inputenc}
\usepackage[russian]{babel}

\usepackage{amsmath,amssymb,amsthm,mathtools}
\usepackage{helvet}
\renewcommand{\familydefault}{\sfdefault}

\usepackage{icomma}
\usepackage[a4paper,margin=22mm,headheight=15pt]{geometry}
\usepackage{microtype}

\usepackage{tikz}
\usetikzlibrary{calc}
\usetikzlibrary{arrows.meta,positioning,shapes.geometric}

\usepackage{tkz-euclide}

\usepackage{pgfplots}
\pgfplotsset{compat=1.18}

\usepackage{tabularx,booktabs,longtable}
\usepackage{enumitem}
\usepackage[hidelinks]{hyperref}
\usepackage{parskip}
\usepackage{fancyhdr}
\usepackage{setspace}
\usepackage{xcolor}

% ------------------------------------------------
% Цветовая палитра
% ------------------------------------------------

\definecolor{accent}{HTML}{008080}
\definecolor{darkaccent}{HTML}{004D4D}
\definecolor{textgray}{HTML}{333333}
\definecolor{softgray}{HTML}{F2F4F4}

% ------------------------------------------------
% Общая типографика
% ------------------------------------------------

\onehalfspacing

\setlength{\parindent}{0pt}
\setlength{\parskip}{0.55em}

\setlist[itemize]{
    leftmargin=*,
    itemsep=0.25em,
    topsep=0.25em
}

\setlist[enumerate]{
    leftmargin=*,
    itemsep=0.45em,
    topsep=0.25em
}

\renewcommand{\arraystretch}{1.28}

% ------------------------------------------------
% Колонтитулы
% ------------------------------------------------

\pagestyle{fancy}
\fancyhf{}

\fancyhead[L]{%
    \small\textcolor{textgray}{Ксения Васильченко, 6 класс}
}

\fancyhead[R]{%
    \small\textcolor{textgray}{Чек-лист}
}

\fancyfoot[C]{%
    \small\thepage
}

\renewcommand{\headrulewidth}{0.3pt}
\renewcommand{\footrulewidth}{0pt}

% ------------------------------------------------
% Заголовки
% ------------------------------------------------

\newcommand{\sectiontitle}[1]{%
    \vspace{0.4em}
    {\Large\bfseries\textcolor{darkaccent}{#1}}\par
    \vspace{0.2em}
    \noindent
    \textcolor{accent}{\rule{\linewidth}{0.7pt}}\par
    \vspace{0.35em}
}

\newcommand{\subhead}[1]{%
    \vspace{0.25em}
    {\large\bfseries\textcolor{darkaccent}{#1}}\par
    \vspace{0.1em}
}

% ------------------------------------------------
% Стили блок-схем
% ------------------------------------------------

\tikzset{
    fcstart/.style={
        ellipse,
        draw=darkaccent,
        line width=0.9pt,
        align=center,
        minimum width=4cm,
        minimum height=1.2cm,
        inner sep=6pt,
        font=\small
    },
    fcproc/.style={
        rectangle,
        rounded corners=4pt,
        draw=darkaccent,
        line width=0.9pt,
        align=center,
        text width=10.8cm,
        minimum height=1.25cm,
        inner sep=7pt,
        font=\small
    },
    fcdecision/.style={
        diamond,
        aspect=2.3,
        draw=darkaccent,
        line width=0.9pt,
        align=center,
        text width=5cm,
        inner sep=2pt,
        font=\small
    },
    fcarrow/.style={
        -{Stealth[length=3mm,width=2mm]},
        draw=darkaccent,
        line width=0.9pt
    },
    fclabel/.style={
        font=\small,
        fill=white,
        inner sep=1.5pt
    }
}

\begin{document}

\color{textgray}

% ================================================================
% Титульный лист
% ================================================================

\begin{titlepage}

\thispagestyle{empty}
\centering

\vspace*{2.8cm}

{\Large\textcolor{accent}{ЧЕК-ЛИСТ}}\par

\vspace{0.8cm}

{\Huge\bfseries\textcolor{darkaccent}{%
НОД, дроби и уравнения%
}}\par

\vspace{0.35cm}

{\Large\bfseries\textcolor{darkaccent}{%
Ключевые идеи занятия%
}}\par

\vspace{1.35cm}

{\large 6 класс}\par

\vspace{0.6cm}

{\large Ксения Васильченко}\par

\vfill

{\large \today}\par

\vspace{0.55cm}

{\large
Лёвин Артём Александрович
эксклюзивно для Ксении Васильченко
}\par

\vspace{1.3cm}

% Сдержанная кривая Безье внизу титульной страницы
\begin{tikzpicture}[remember picture,overlay]

\draw[
    accent,
    line width=1.4pt
]
([xshift=16mm,yshift=14mm]current page.south west)
..
controls
([xshift=58mm,yshift=27mm]current page.south west)
and
([xshift=-72mm,yshift=6mm]current page.south east)
..
([xshift=-16mm,yshift=18mm]current page.south east);

\end{tikzpicture}

\end{titlepage}

% ================================================================
% Основной чек-лист
% ================================================================

\sectiontitle{Основной чек-лист}

% ------------------------------------------------
% Десятичные дроби
% ------------------------------------------------

\subhead{1. Десятичные дроби: быстрый контроль вычислений}

\begin{itemize}

\item
При сложении и вычитании выравнивайте числа
\textbf{по запятой}.

При необходимости дописывайте нули:
\[
1{,}3=1{,}30.
\]

\item
При умножении сначала умножайте числа как натуральные,
затем отделяйте справа столько десятичных знаков,
сколько их было \textbf{суммарно} в обоих множителях.

Например:
\[
0{,}24\cdot1{,}3=0{,}312.
\]

\item
При делении десятичной дроби на натуральное число
следите за положением запятой в частном:
\[
0{,}24:4=0{,}06.
\]

\end{itemize}

% ------------------------------------------------
% Обыкновенные дроби
% ------------------------------------------------

\subhead{2. Обыкновенные и смешанные дроби}

\begin{itemize}

\item
Для сложения и вычитания дробей нужен
\textbf{общий знаменатель}.

Например:
\[
\frac23+\frac15
=
\frac{10}{15}+\frac{3}{15}
=
\frac{13}{15}.
\]

\item
При приведении дроби к новому знаменателю
числитель и знаменатель умножаются
на одно и то же число.

\item
Смешанное число удобно переводить
в неправильную дробь:
\[
2\frac23
=
\frac{2\cdot3+2}{3}
=
\frac83.
\]

\item
При умножении дробей перемножайте числители
и знаменатели:
\[
\frac23\cdot\frac75
=
\frac{14}{15}.
\]

\item
Перед умножением дробей полезно выполнять сокращение,
если числитель и знаменатель содержат общий множитель.

Сокращаются именно \textbf{множители}.

\item
Деление на дробь заменяется умножением
на обратную дробь:
\[
\frac23:\frac75
=
\frac23\cdot\frac57
=
\frac{10}{21}.
\]

\item
Неправильную дробь можно перевести
в смешанную делением числителя на знаменатель:
\[
\frac{16}{15}
=
1\frac{1}{15}.
\]

\end{itemize}

\textbf{Проверочный пример}

\[
2\frac23-1\frac35
=
\frac83-\frac85
=
\frac{40}{15}-\frac{24}{15}
=
\frac{16}{15}
=
1\frac1{15}.
\]

% ------------------------------------------------
% НОД
% ------------------------------------------------

\subhead{3. Разложение на простые множители и НОД}

\textbf{Наибольший общий делитель} двух натуральных чисел ---
это наибольшее натуральное число,
на которое оба числа делятся без остатка.

При разложении на простые множители
для НОД берутся только
\textbf{общие простые множители}
в наименьших степенях.

Например:
\[
84=2^2\cdot3\cdot7,
\qquad
112=2^4\cdot7.
\]

Общие множители:
\[
2^2\cdot7.
\]

Следовательно,
\[
\text{НОД}(84,112)=28.
\]

\textbf{Самопроверка:}
\[
84:28=3,
\qquad
112:28=4.
\]

Оба частных являются натуральными числами.

% ------------------------------------------------
% Смысл НОД
% ------------------------------------------------

\subhead{4. Как понять смысл НОД в текстовой задаче}

НОД особенно полезен,
когда несколько количеств нужно разбить
на максимально возможные одинаковые части,
группы, наборы или упаковки
без остатка.

Главный вопрос:

\begin{center}
\textbf{Какая величина по условию должна быть одинаковой
и максимально возможной?}
\end{center}

\begin{itemize}

\item
Если требуется максимальное
\textbf{число одинаковых наборов},
то НОД показывает число наборов.

Деление исходных количеств на НОД
показывает состав одного набора.

\item
Если требуется максимальное
\textbf{число объектов в каждой группе},
то НОД показывает размер одной группы.

Деление исходных количеств на НОД
показывает количество групп.

\end{itemize}

\textbf{Пример 1}

Из $84$ помидоров и $112$ огурцов
делают максимальное число одинаковых наборов,
используя все овощи.

\[
\text{НОД}(84,112)=28.
\]

Следовательно, получится
\[
28
\]
одинаковых наборов.

Состав каждого:
\[
84:28=3
\]
помидора и
\[
112:28=4
\]
огурца.

\textbf{Пример 2}

$108$ мальчиков и $114$ девочек
распределяют отдельно по группам
с одинаковым максимально возможным
числом человек в каждой группе.

\[
108=2^2\cdot3^3,
\]

\[
114=2\cdot3\cdot19.
\]

Поэтому
\[
\text{НОД}(108,114)=2\cdot3=6.
\]

Значит, в каждой группе находится
по $6$ человек.

Количество групп мальчиков:
\[
108:6=18.
\]

Количество групп девочек:
\[
114:6=19.
\]

\textbf{Типичная ошибка:}
считать, что НОД всегда означает количество групп.

Смысл найденного НОД определяется
формулировкой конкретной задачи.


% ================================================================
% Уравнения
% ================================================================

\sectiontitle{Уравнения}

\subhead{5. Связь компонентов арифметических действий}

Удобно понимать смысл арифметического действия
и проверять правило простым числовым примером.

\begin{table}[ht]

\centering

\begin{tabularx}{\linewidth}{
    @{}
    >{\bfseries}l
    X
    X
    @{}
}

\toprule

Равенство
&
Как найти первый неизвестный компонент
&
Как найти второй неизвестный компонент
\\

\midrule

$a+b=c$
&
$a=c-b$
&
$b=c-a$
\\

$a-b=c$
&
$a=b+c$
&
$b=a-c$
\\

$ab=c$
&
$a=\dfrac{c}{b}$,
если $b\ne0$
&
$b=\dfrac{c}{a}$,
если $a\ne0$
\\

$\dfrac{a}{b}=c$
&
$a=bc$,
причём $b\ne0$
&
$b=\dfrac{a}{c}$,
если $c\ne0$;
обязательно $b\ne0$
\\

\bottomrule

\end{tabularx}

\end{table}

\textbf{Удобная модель}

Для равенства
\[
a-b=c
\]
можно взять числовой пример
\[
10-2=8.
\]

Чтобы найти уменьшаемое $10$:
\[
8+2=10.
\]

Чтобы найти вычитаемое $2$:
\[
10-8=2.
\]

Такой числовой пример помогает проверить правило,
не опираясь только на запоминание формулы.

% ------------------------------------------------
% Простые уравнения
% ------------------------------------------------

\subhead{6. Простые уравнения}

\textbf{Пример 1}
\[
\frac{x}{2}=10.
\]

Чтобы найти делимое:
\[
x=10\cdot2=20.
\]

Проверка:
\[
20:2=10.
\]

\textbf{Пример 2}
\[
3x=15.
\]

Тогда
\[
x=15:3=5.
\]

\textbf{Пример 3}
\[
x-3=10.
\]

Тогда
\[
x=10+3=13.
\]

\textbf{Пример 4}
\[
15-x=10.
\]

Чтобы найти вычитаемое:
\[
x=15-10=5.
\]

% ------------------------------------------------
% Неизвестное в знаменателе
% ------------------------------------------------

\subhead{7. Неизвестное в знаменателе}

Если выражение с неизвестной находится
в знаменателе, знаменатель
\textbf{не может быть равен нулю}.

Например:
\[
\frac{2}{x}=10.
\]

Ограничение:
\[
x\ne0.
\]

Для равенства
\[
\frac{a}{b}=c
\]
делитель можно найти как
\[
b=\frac{a}{c}.
\]

Поэтому
\[
x=\frac{2}{10}
=\frac15.
\]

Проверка:
\[
\frac{2}{1/5}=10.
\]

% ------------------------------------------------
% Составные уравнения
% ------------------------------------------------

\subhead{8. Составные уравнения}

Главный приём:
временно рассматривайте целое выражение,
содержащее $x$,
как один неизвестный компонент.

\textbf{Пример 1}

\[
\frac{x+3}{2}=10.
\]

Выражение $x+3$ играет роль делимого.

Поэтому
\[
x+3=10\cdot2=20.
\]

Далее:
\[
x=20-3=17.
\]

Проверка:
\[
\frac{17+3}{2}=10.
\]

\textbf{Пример 2}

\[
\frac{6}{x-5}=2.
\]

Сначала ограничение:
\[
x-5\ne0,
\qquad
x\ne5.
\]

Выражение $x-5$ играет роль делителя:
\[
x-5=\frac62=3.
\]

Поэтому
\[
x=3+5=8.
\]

Проверка:
\[
\frac{6}{8-5}=2.
\]

Условие
\[
x\ne5
\]
также выполнено.

\textbf{Пример 3}

\[
\frac{14-x}{2}=5.
\]

Сначала:
\[
14-x=10.
\]

Теперь нужно найти вычитаемое:
\[
x=14-10=4.
\]

Проверка:
\[
\frac{14-4}{2}=5.
\]

% ------------------------------------------------
% Типичные ошибки
% ------------------------------------------------

\subhead{9. Типичные ошибки}

\begin{itemize}

\item
Не меняйте знаки механически.
Каждый переход должен соответствовать
конкретному обратному действию.

\item
Если неизвестное находится в знаменателе,
сначала запишите ограничение:
знаменатель не равен нулю.

\item
После решения подставьте найденное значение
в исходное уравнение.

\item
При сокращении дробей сокращайте
общие \textbf{множители}.
Отдельные слагаемые суммы сокращать нельзя.

\item
При сложении и вычитании дробей
сначала приводите их к общему знаменателю.

\item
При работе со смешанными числами
следите за правильным переводом
в неправильную дробь.

\item
В текстовой задаче на НОД
обязательно подпишите,
что означает найденное число.

\end{itemize}


% ================================================================
% Тренировочный блок
% ================================================================

\sectiontitle{Тренировочный блок --- Путь А}

\textbf{Выполните 5 заданий.}

Решения оформляйте аккуратно.
В текстовых задачах обязательно указывайте,
что означает каждое найденное число.

\begin{enumerate}

% -------------------- 1 --------------------

\item
Выполните вычисления:

\[
\text{а) }\;
1{,}7-0{,}46;
\]

\[
\text{б) }\;
0{,}32\cdot1{,}5;
\]

\[
\text{в) }\;
\frac34+\frac25;
\]

\[
\text{г) }\;
2\frac13-1\frac34.
\]

% -------------------- 2 --------------------

\item
Разложите числа
\[
72
\quad\text{и}\quad
120
\]
на простые множители.

Найдите
\[
\text{НОД}(72,120).
\]

% -------------------- 3 --------------------

\item
Есть $96$ яблок и $144$ груши.

Из всех фруктов нужно составить
максимальное число одинаковых наборов,
использовав все яблоки и все груши.

Определите:

\begin{itemize}

\item
сколько одинаковых наборов получится;

\item
сколько яблок будет в каждом наборе;

\item
сколько груш будет в каждом наборе.

\end{itemize}

% -------------------- 4 --------------------

\item
$126$ мальчиков и $168$ девочек
распределяют отдельно по группам.

Во всех группах должно быть
одинаковое максимально возможное
число человек.

Определите:

\begin{itemize}

\item
сколько человек будет в одной группе;

\item
сколько получится групп мальчиков;

\item
сколько получится групп девочек.

\end{itemize}

% -------------------- 5 --------------------

\item
Решите уравнения.
Для каждого уравнения выполните проверку.

\[
\text{а) }\;
\frac{x+7}{3}=9;
\]

\[
\text{б) }\;
\frac{12}{x-2}=3.
\]

Для пункта б) сначала укажите ограничение
на значение $x$.

\end{enumerate}

% ================================================================
% Ответы
% ================================================================

\clearpage

\sectiontitle{Ответы}

\begin{table}[ht]

\centering

\begin{tabularx}{\linewidth}{
    @{}
    c
    X
    @{}
}

\toprule

\textbf{№}
&
\textbf{Ответ}
\\

\midrule

1
&
а) $1{,}24$;

б) $0{,}48$;

в) $\dfrac{23}{20}=1\dfrac3{20}$;

г) $\dfrac7{12}$.
\\[1mm]

2
&
\[
\text{НОД}(72,120)=24.
\]
\\

3
&
$48$ одинаковых наборов;

в каждом наборе по $2$ яблока
и $3$ груши.
\\

4
&
По $42$ человека в каждой группе;

$3$ группы мальчиков;

$4$ группы девочек.
\\

5
&
а) $x=20$;

б) $x=6$,
при ограничении
\[
x\ne2.
\]
\\

\bottomrule

\end{tabularx}

\end{table}

\vfill

\begin{center}

\small
\textcolor{textgray}{%
Лёвин Артём Александрович
эксклюзивно для Ксении Васильченко%
}

\end{center}

\end{document}